% =====================================================================
\chapter{Literature Map and Positioning}\label{ch:literature}
% =====================================================================

This chapter places the paper's claims in their literature context, states their relation
to existing theory, and marks where the gaps lie.

\section{Six Groups of Literature}\label{sec:six-groups}

\paragraph{First Group: Modern Specification Generation and Joint Evaluation.}
This is the frontier closest to the intent $\to$ specification step. The relevant
literature reports measurable results: a specification-generation benchmark for
operating-system kernel verification tasks still shows model limits on long context and
specification completeness \cite{osvbench2026}; a benchmark that evaluates code
correctness, specification quality and proof generation separately reports a significant
gap among the three dimensions (on a particular task set, code correctness, specification
soundness and completeness, and single-pass proof success rate are 72.6\%, 52.3\% and
4.9\% respectively) \cite{verina2026}; an evaluation framework that instantiates
specifications as trusted tests points out that proof difficulty interferes with the
assessment of specification quality \cite{coins2026}; a traceable refinement method
decomposes natural-language requirements into atomic requirements and establishes
requirement-level traceability \cite{verispecgen2026}; a mutation-driven method tests
specification strength with semantically non-equivalent variants \cite{specsyn2026}; and
the route of generating specifications from code should be studied separately from the
route of generating specifications from intent \cite{specgen2025}. Vacuity detection
\cite{kupferman2003} supplies a theoretical route to identifying the case in which
verification passes while the specification holds vacuously. Together these results support
one important judgement: \emph{code generation ability, specification expressiveness and
formal proof ability are distinct research variables and should be evaluated separately}.

\paragraph{Second Group: Program Semantics, Refinement and Observable Behaviour.}
Operational semantics, denotational semantics, the lattice model of abstract interpretation
\cite{cousot1977}, monadic semantics \cite{moggi1991,plotkin2002}, domain theory,
refinement relations \cite{hoare1969} and separation logic \cite{appel2007} together supply
the mathematical tools for the question of when two systems at different levels implement
the same behavioural specification. Semantic refinement is not simple information
compression: it must define the \emph{behavioural relation} between levels.

\paragraph{Third Group: Type Theory, Constructive Mathematics and Proof Extraction.}
Constructive type theory \cite{hott2013}, the correspondence between proofs and programs
\cite{howard1980}, and the mechanisms that connect proof objects with constructive
computation supply the mathematical foundation for ``an AI proposes candidate proofs and a
mechanical kernel checks them''. One inference must be avoided: \emph{if a specification
has already been translated into a type, then the specification necessarily expresses human
intent faithfully}. The correspondence guarantees only the stipulated structural
correspondence among propositions, types, proofs and terms within a chosen logic and type
system; how human intent is turned into a proposition is a separate process that requires
verification. Refinement types \cite{rondon2008} and SMT solving \cite{z32008} are the
engineering representatives of this direction; automated theorem-proving environments and
retrieval-augmented proof search \cite{leandojo2023,leancopilot2024} are its continuation
in the age of AI.

\paragraph{Fourth Group: From Proof to Machine Code.}
This is the most mature group. A verified compiler turns the compilation transformation
itself into a proved object \cite{leroy2009}; the back ends of a verified compiler cover
several target architectures \cite{cakeml2019}; a refinement proof for an operating-system
kernel shows how far the formal verification of large-scale low-level system software can
reach \cite{klein2009}; translation validation checks a \emph{single} compilation result
independently \cite{necula2000}; translation certification establishes a verifiable
relational certificate for a \emph{specific artefact} \cite{krijnen2024}; and proof-carrying
code lets a code provider attach a proof that a machine can check independently
\cite{necula1997}.

The key lesson of this group is that \emph{verifying the correctness of a compiler and
verifying the correctness of a particular binary are two different strategies}. The former
proves that the compiler satisfies a universal property; the latter verifies one build
output independently. The two can be combined, without staking all trust on a single tool.

\paragraph{Fifth Group: Resource Correctness.}
Automatic amortised resource analysis can derive symbolic resource bounds for a program
under a particular language and cost model, with a soundness proof; worst-case execution
time analysis derives timing bounds for particular hardware and assumptions; static cache
analysis shows the influence of the hardware model on bound inference. Three notions of
time must be distinguished: algorithmic complexity gives a bound on the growth of cost as
input size varies; worst-case execution time analysis gives an upper bound under a
particular platform and assumptions; measurement provides latency data under particular
experimental conditions. \emph{None of the three substitutes for another}. Equally, linear
use is not optimal use: a program may consume each input resource exactly once and still use
a poor algorithm or introduce excessive latency.

\paragraph{Sixth Group: Runtime Verification, Contract Composition and the Emergence of Levels.}
This group contains two classes of material of different kinds, which must be treated
separately.

The first class is \emph{runtime assurance and contracts}: assume--guarantee reasoning
requires a subsystem to give the corresponding guarantee when its input assumption is
satisfied \cite{henzinger1998}; contract composition for dynamical systems gives composition
rules and checking methods under arbitrary interconnection topologies
\cite{sharf2021,sharf2024,shali2026}; work on the uncertainty of runtime verification shows
that observation, event transport and trace integrity are components of monitoring
reliability \cite{taleb2023}; online monitoring of temporal logic \cite{leucker2009} and
information-flow security \cite{goguen1982} address trace properties and access
restrictions respectively.

The second class is \emph{explanatory frameworks}: the near-decomposability of complex
systems provides an explanation for level structures \cite{simon1962}; minimum description
length supplies the statistical and information-theoretic background for ``abstraction is
compression'' \cite{rissanen1978,grunwald2006,grunwald2019}; computational thermodynamics
gives a physical lower bound on computation \cite{landauer1961}. This group can supply an
explanatory framework for architecture, compression, model selection and complex systems,
but it \emph{cannot automatically supply a correctness proof for a particular program}.

\section{The Five Layers Rearranged by Mathematical Function}\label{sec:five-layers}

Rearranged by which segment of the chain it solves, this literature yields a structure
deeper than ``an AI writes code and a verifier checks it'':

\begin{table}[htbp]
\centering
\small
\caption{The five layers rearranged by mathematical function}\label{tab:five-layers}
\begin{tabular}{L{1.6cm}L{5.4cm}L{7.2cm}}
\toprule
\textbf{Layer} & \textbf{Function} & \textbf{Representative lines of work} \\
\midrule
A & Expressibility of the objective: formal specification generation, intent traceability,
  requirement decomposition, counterexample generation &
  Specification-generation benchmarks, traceable refinement, mutation-driven strength
  testing, vacuity detection, requirement formalisation \\
B & Realisability of the objective: logical consistency, program synthesis, constraint
  solving, realisability decision &
  Counterexample-guided inductive synthesis, syntax-guided synthesis, program-skeleton
  completion, constructive type theory, SMT solving \\
C & Cross-layer semantic preservation: operational semantics, refinement relations, proof
  checking, compiler correctness, binary verification &
  Verified compilers, verified compiler back ends, kernel refinement proofs, proof-carrying
  code, translation validation \\
D & Resource and environment assurance: complexity, resource bounds, runtime monitoring,
  assume--guarantee contracts, execution constraints &
  Automatic amortised resource analysis, worst-case execution time analysis, runtime
  verification, contract composition, information-flow security \\
E & Structure and physical limits: model compression, levels of complex systems,
  thermodynamics and physical realisation &
  Minimum description length, near-decomposability of complex systems, computational
  thermodynamics, molecular self-assembly \\
\bottomrule
\end{tabular}
\end{table}

The most important finding is that these areas \emph{do not share one ready-made, fully
unified mathematical object}. Each has nonetheless solved a different problem along this
chain. The research opportunity therefore lies in defining the composition relations among
them, not in inventing the entire mathematical toolkit from scratch.

\section{The Unoccupied Intersection}\label{sec:gaps}

\begin{principle}[Shape of the literature map]\label{pr:map-shape}
Every step from specification to binary has a mature literature; automated synthesis from
proof to implementation has an active frontier; but the gap from \emph{intent} to
\emph{specification} has no formal solution, only consensus structures, and the literatures
of social epistemology, Byzantine fault tolerance and adversarial review exist in entirely
different disciplines and have not been systematically coupled with formal verification
systems.
\end{principle}

On these grounds, the most valuable position to occupy is the following.

\begin{principle}[Proposed research position]\label{pr:position}
Integrate adversarial consensus in the style of social epistemology with mechanical
admission in the style of proof-carrying code inside one framework: let a
consensus-attested specification be the direct input to a proof-carrying artefact, and
anchor consensus records, specification versions and artefact identity together
cryptographically. As far as the review conducted for this paper extends, few works address
both literatures at once.
\end{principle}

This position is valuable because it satisfies three conditions at once: it corresponds to a
real and still unsolved technical gap (no mechanical solution at the intent layer); both of
its components already have mature theoretical foundations; and its composition rules can
be written out explicitly (Definition \ref{def:two-layer} and Theorem \ref{th:assurance})
and can therefore be tested.

\section{Relation to Existing Theory}\label{sec:relations-lit}

\begin{longtable}{L{3.6cm}L{10.8cm}}
\caption{Relation of this paper to existing theory}\label{tab:relations}\\
\toprule
\textbf{Theory} & \textbf{How this paper uses it, and how it does not use it} \\
\midrule
\endfirsthead
\toprule
\textbf{Theory} & \textbf{How this paper uses it, and how it does not use it} \\
\midrule
\endhead
\bottomrule
\endlastfoot
Near-decomposability of complex systems &
  \emph{Used}: as the explanation of why level structures are stable, and recast into a
  testable sufficient-statistic condition (Definition \ref{def:neardecomp}, Proposition
  \ref{prop:leakpair}).
  \emph{Not used}: as evidence that ``all complex systems necessarily have the same
  levels''; the source work does not prove this, nor does it give a general theorem about
  verification cost. \\
Abstract interpretation and Galois connections &
  \emph{Used}: as the formal tool for semantic preservation between levels (Definition
  \ref{def:galois}).
  \emph{Not used}: no claim that all inter-level relations can be expressed as Galois
  connections; where they cannot, the paper falls back on simulation relations or on an
  explicit human argument. \\
Minimum description length &
  \emph{Used}: as the formal background for ``comprehension is compression'' and the source
  of the stopping rule for descent (Definition \ref{def:mdl}, Proposition
  \ref{prop:stopping}).
  \emph{Not used}: no claim that the shortest description is the most correct description.
  MDL is a model-selection principle; its conclusions depend on the chosen model family,
  encoding and data, and it cannot by itself determine what purpose a system should
  achieve. \\
Information hiding and deep modules &
  \emph{Used}: as the normative basis for module boundaries.
  \emph{Not used} as it stands: the paper adds a testable version of it, namely whether an
  interface is a sufficient statistic of behaviour (Definition \ref{def:neardecomp}),
  because how well information is hidden is not itself measurable. \\
Correspondence between organisational structure and system structure &
  \emph{Used}: as one of the mechanisms of high-level convergence and one of the causes of
  divergence.
  \emph{Not used}: not as a technical necessity; it is a source of constraints, not a
  source of uniqueness for solutions. \\
Program logic and refinement &
  \emph{Used}: as the mathematical interface for the specification--implementation
  relation.
  \emph{Not used}: no claim that safety refinement automatically guarantees liveness,
  timing and fairness properties. \\
Proof-carrying code &
  \emph{Used}: as the classic implementation of a trust structure in which generation is
  free and admission is strict, and as the predecessor of the fourth pillar of this paper.
  \emph{Not used}: its premise is not ignored: the receiver must fix a safety policy in
  advance, and the correctness of the policy itself is an external problem. \\
Assume--guarantee reasoning &
  \emph{Used}: as the source of the rules by which local guarantees compose into a global
  guarantee (Proposition \ref{prop:soundness}).
  \emph{Not used}: its four composition conditions are not ignored; a component's being
  ``correct in isolation'' and a system's being ``correct as a whole'' are not
  interchangeable. \\
Resource analysis &
  \emph{Used}: as evidence that resource correctness can enter formal analysis.
  \emph{Not used}: a symbolic bound is not equated with a time bound on a particular
  processor; complexity order, worst-case execution time, statistical quantiles and
  measurement are not conflated. \\
Computational thermodynamics &
  \emph{Used}: as the basis for the boundary condition that the physical cost of computation
  is incompressible.
  \emph{Not used}: it is not used to predict the power consumption of a piece of ordinary
  code, nor folded into semantic correctness.
  (The paper mentions it only in Chapter \ref{ch:objections}, and not as a load-bearing
  structure.) \\
The halting problem and Rice's theorem &
  \emph{Used}: as the boundary at which a general synthesis algorithm and a general decision
  algorithm cannot exist.
  \emph{Not used}: they are not used to prove that ``intent can never be formalised''
  (Proposition \ref{prop:no-godel}). \\
$\mathrm{P}$ and $\mathrm{NP}$ &
  \emph{Used}: only the structural asymmetry entailed by the \emph{definition} of the class
  $\mathrm{NP}$ (Proposition \ref{prop:asym-struct}).
  \emph{Not used}: no open assertion of ``$\mathrm{P} \neq \mathrm{NP}$'' or
  ``$\mathrm{P} = \mathrm{NP}$'' serves as a load-bearing structure. \\
\bottomrule
\end{longtable}

\begin{principle}[Positioning of this paper]\label{pr:self-position}
The contribution of this paper lies at the \emph{interface layer}, not in a new semantic
theory: it promotes levels from a comprehension aid to the unit into which verification
obligations decompose; it formalises oblivious sets, near-decomposability and read-sets as
computable or measurable objects; it attributes the necessity of internal
infrastructuralisation to the completeness of the operational semantics rather than to code
volume; and it writes the trust structure as the composite of the two modes of trust
generation and gives a bound on consensus reliability. These interfaces connect existing
theory rather than replace it.
\end{principle}
